\subsection{倒班与加班：作息错位下的性生活}

上一节讨论的是"你偏好清晨还是夜晚"——但有一类人的作息不是偏好决定的，而是工作决定的：夜班、轮班、长期加班的人，性生活遇到的问题是系统性的，值得单独说清楚。

\textbf{先看倒班怎么影响身体}。睾酮的分泌高度依赖睡眠，尤其是凌晨的快速眼动睡眠期——夜班和轮班把睡眠切碎、颠倒，直接压低了睾酮水平；同时生物钟紊乱让皮质醇节律异常、褪黑素分泌延迟。多项针对轮班工人的研究得到了一致的结论：\textbf{性欲与性功能的下降在倒班人群中确实更常见}。这不是"身体不行了"，而是身体在错误的时段被迫工作。

再看清一个现实：\textbf{夫妻的"清醒时间窗"几乎不重叠}。一个白天上班、一个夜里上班，两个人的交集只剩换班时疲惫的几小时——性在这样的日子里自然就消失了。认识到这是排期问题而不是感情问题，是解决它的第一步。

可以做的事：

\begin{itemize}
	\item \textbf{把共同的休息日固定为优先窗口}：与其在错位的日常里等"撞上"，不如在日历上把共同休息日的某个时段固定下来——如同本书反复说的"预约性爱"，对倒班家庭它不是浪漫的敌人，而是唯一的现实解；
	\item \textbf{下夜班先补觉，不给性打卡}：下了夜班强行"配合对方的时间"往往两头失败。先睡够，在精力最好的那几个小时里亲密，质量远胜于硬撑；
	\item \textbf{错位期不等于零亲密}：见不到面的时段，一句语音、一段留言、出门前认真的拥抱，都是在为下次见面"充电"；
	\item \textbf{轮换排班时提前通气}：如果换班时间表可以预知，把"下一次共同休息"提前告知彼此——有盼头的等待比无尽的错位好得多。
\end{itemize}

\textbf{长期加班的问题与倒班类似，但更隐蔽}：没有夜班那么极端，却是"每天都累一点"。慢性疲惫意味着身体长期把性欲排在最低优先级——持续偏高的皮质醇会抑制下丘脑—垂体—性腺轴，性欲下降是生理后果，不是态度问题。对这类伴侣，"低耗能亲密"比"恢复出厂设置"更现实：不追求完整流程的依偎、按摩、一起泡脚、睡前十分钟的聊天，先恢复"身体接触是安全的、不需要回报的"这个感受，性欲往往在这种松动里自己回来。\textbf{警惕把性变成又一项 KPI}——加班家庭最不需要的就是床上的考核。

一条界线：如果性欲的消失是长期的（几个月以上），并且伴随着情绪低落、对大多数事情提不起兴趣、睡眠差得超出工作所能解释的范围——请把抑郁与甲状腺问题列入排查清单，和伴侣一起去见一次医生。另外，因工作长期分居两地的伴侣（出差、异地），可参见下一节的出行提示与数字亲密相关章节；重逢之夜的预期管理，本书在异地重逢相关内容中已有专述。
